% Source: https://artofproblemsolving.com/community/c3h1347915_cyclic_hexagon
% Stable identity: aops-topic-1347915-q01
Let $R$ be the radius and $O$ the centre. A chord of length $\ell$ subtends a central angle $2\theta$ with $\sin\theta=\dfrac{\ell}{2R}$. Let the sides of length $2$ subtend $2\theta$ and the sides of length $12$ subtend $2\varphi$, so
\[
\sin\theta=\dfrac{1}{R},\qquad \sin\varphi=\dfrac{6}{R} .
\]

\begin{asy}
usepackage("amsmath");
size(6cm);
defaultpen(fontsize(10pt));
real R=9; real al=2*aSin(1/9), be=2*aSin(2/3);
pair[] P; real th=90+al/2;
real[] st={al,be,be,al,be,be};
for(int i=0;i<6;++i){P.push(R*dir(th)); th-=st[i];}
pair O=(0,0);
draw(circle(O,R),gray);
draw(P[0]--P[1]--P[2]--P[3]--P[4]--P[5]--cycle,linewidth(0.9));
draw(O--P[0],blue+dashed); draw(O--P[1],blue+dashed); draw(O--P[2],red+dashed);
dot(O);
label("$O$",O,S);
for(int i=0;i<6;++i){pair m=(P[i]+P[(i+1)%6])/2; label(i%3==0 ? "$2$" : "$12$",m,unit(m));}
label("$R$",(O+P[2])/2,S,red);
label("\mbox{}",(min(currentpicture,user=true)+max(currentpicture,user=true))/2,black); // keeps the text after the figure black when LaTeX typesets the labels inline
\end{asy}

The central angles add up to $360^\circ$: $2\cdot2\theta+4\cdot2\varphi=360^\circ$, so $\theta=90^\circ-2\varphi$. Then
\[
\sin\theta=\cos2\varphi=1-2\sin^2\varphi\quad\Longrightarrow\quad \dfrac{1}{R}=1-\dfrac{72}{R^2}\quad\Longrightarrow\quad R^2-R-72=0 ,
\]
so $(R-9)(R+8)=0$ and $R=9$. (Check: $\sin\theta=\dfrac{1}{9}$, $\sin\varphi=\dfrac{2}{3}$, and $\cos2\varphi=1-\dfrac{8}{9}=\dfrac{1}{9}$.) The area of the circle is
\[
\pi R^2=\boxed{81\pi}.
\]
